\section{尺度分离与渐近近似}
\label{sec:15-Differential-Equations-and-Dynamical-Systems/01-ODE-Theory/04}

一个化学反应网络的仿真跑得出奇地慢。你查了半天，发现积分器把步长压到了 \(10^{-7}\) 秒，而你关心的浓度变化发生在几十分钟的尺度上。中间隔着十个数量级，机器却必须把每一微秒都算一遍。

这类模型里通常有几个反应快得几乎瞬间完成。它们的快，恰恰是可以拿来简化问题的东西：既然快变量在慢变量还没动之前就已经落到平衡，那就直接把它当成平衡，从模型里消掉。本节讲的是怎样把这个直觉做成可控的近似，以及它在哪里会骗人。

\subsection{先无量纲化，才知道什么算小}

长度 \(x\)、时间 \(t\) 和温度差 \(T\) 带着单位时，直接比较方程中两项的数值大小，比出来的可能只是单位选择。取典型尺度 \(L,\tau,\Delta T\)，令

\[
\widehat x=\frac{x}{L},
\qquad
\widehat t=\frac{t}{\tau},
\qquad
\widehat T=\frac{T}{\Delta T},
\]

方程变成无量纲形式，材料性质、速度与尺度会自动组合成少数几个无量纲参数。只有这样得到的 \(\varepsilon\ll1\) 才有稳定含义：它表示两个物理尺度之比很小，而不是某个带单位的数字碰巧写得小。换一套单位，前者不变，后者面目全非。

\subsection{快变量落到慢流形上就不再是自由度}

尺度分离最常用的形态是这样一个系统：

\[
x'=f(x,y),
\qquad
\varepsilon\,y'=g(x,y),
\]

其中 \(x\) 慢、\(y\) 快，\(\varepsilon\ll1\)。第二个方程写成 \(y'=g(x,y)/\varepsilon\)，右端被放大了 \(1/\varepsilon\) 倍：只要 \(g\) 不为零，\(y\) 就以极快的速度冲向使 \(g=0\) 的那个值。

\begin{center}
\begin{tikzpicture}[>=stealth,xscale=1.55,yscale=1.5]
  \draw[->] (-0.1,0) -- (4.5,0) node[right] {\(x\)（慢）};
  \draw[->] (0,-0.1) -- (0,2.9) node[above] {\(y\)（快）};
  \draw[thick] plot[smooth] coordinates
    {(0,0.50) (0.5,0.86) (1,1.15) (1.5,1.36) (2,1.50) (2.5,1.56) (3,1.55) (3.5,1.46) (4,1.30)};
  \node[right] at (4.02,1.28) {\small \(g(x,y)=0\)};
  \foreach \xx/\gg in {0.7/0.98, 1.7/1.42, 2.7/1.56, 3.7/1.40}{
    \draw[gray,->] (\xx,{\gg+1.05}) -- (\xx,{\gg+0.14});
    \draw[gray,->] (\xx,{\gg-1.05}) -- (\xx,{\gg-0.14});
  }
  \draw[very thick,->] plot[smooth] coordinates
    {(0.30,2.60) (0.33,1.90) (0.37,1.25) (0.46,0.92) (0.60,0.93) (1,1.15) (1.5,1.36) (2,1.50) (2.5,1.56) (3,1.55)};
  \node[right] at (0.36,2.30) {\small 快过程};
  \node[below] at (2.0,1.30) {\small 慢过程};
\end{tikzpicture}
\end{center}

\emph{快变量先近乎垂直地落到慢流形 \(g=0\) 上，之后整段演化被慢变量牵着沿曲线爬行。}

图上的灰色箭头是快场：离开曲线，\(y\) 几乎沿垂直方向被拽回去。粗线是一条真实轨迹，它在极短的时间里完成第一段下落，此后就贴着曲线走，纵坐标的变化完全由横坐标带动。这段贴合之后，\(y\) 不再是独立的自由度——解 \(g(x,y)=0\) 得到 \(y=h(x)\)，代回第一个方程，模型降维成

\[
x'=f\bigl(x,h(x)\bigr).
\]

这就是把快变量当作瞬时平衡的意思，化学里叫准稳态假设。开头那个反应网络之所以逼得步长降到微秒，是因为积分器必须分辨那段垂直下落；换成降维后的慢模型，需要分辨的只剩下沿曲线的爬行。

代价写在图上：第一段下落被整个抹掉了。若你关心的恰是快过程本身，或者初值离曲线很远而下落途中会撞上某个约束，降维模型不会告诉你。曲线的折返处（图中隆起的顶点附近）也危险，那里 \(g=0\) 可能失去唯一解，快变量会跳到另一支上，慢流形近似在那一点断裂。

这件事与数值分析里的刚性是同一现象的两种处理。\hyperref[sec:08-Numerical-Analysis/06-ODE/03]{``稳定区域与刚性方程''}中，刚性的定义正是系统里同时存在跨若干数量级的时间常数，显式方法的步长被最快的那个模态卡死，即便它早已衰减到看不见。降维是从模型侧消掉快模态，隐式方法是从算法侧容忍它——两条路解决的是同一个 \(\varepsilon\)。

\subsection{正则摄动允许逐阶修正}

尺度分离的另一种用法是把小参数当展开变量。考虑

\[
y'(t)+y(t)=\varepsilon y(t)^2,
\qquad y(0)=a,
\]

并设

\[
y(t;\varepsilon)=y_0(t)+\varepsilon y_1(t)+\varepsilon^2y_2(t)+\cdots.
\]

代回方程按 \(\varepsilon\) 的幂比较系数，依次得到

\[
y_0'+y_0=0,
\qquad
y_1'+y_1=y_0^2.
\]

非线性问题被拆成一串线性问题：首项给主要行为，后续项逐级修正，每一级的右端都由前面已经解出的东西提供。这样的展开若在所关心的时间区间上误差一致受控，称为正则摄动。它给的是带适用范围的近似；把无穷级数形式写出来，并不自动等于得到了收敛解。

\subsection{小参数乘在最高阶导数上，极限就不再均匀}

危险情形是小参数压在最高阶导数前面。例如

\[
\varepsilon y''+y'=0,
\qquad
y(0)=0,\quad y(1)=1.
\]

直接令 \(\varepsilon=0\)，二阶方程退化成一阶，只能满足一个边界条件。方程的阶数变了，说明这个极限不均匀：在区间的大部分地方，解接近缓慢变化的外部解；而在某个边界附近宽度 \(O(\varepsilon)\) 的窄带里，导数急剧增大，用来迅速凑满另一个边界条件。

引入拉伸变量 \(\xi=x/\varepsilon\) 可以把这条窄带放大到 \(O(1)\) 尺度，得到内部解，再让内部解与外部解在共同有效的重叠区对上，就是匹配渐近展开。结构上它与上一小节那张图同源：窄带对应快过程，外部解对应沿慢流形的爬行，区别只在于快慢的分界这次落在空间边界而不是时间起点。薄涂层中的急剧温降、车辆启动瞬间的机械响应与电子控制的慢变化，都是这种快层加慢主体的样子。

\subsection{近似必须说明极限方式和观察窗口}

\(O(\varepsilon)\) 的严格含义是：存在与足够小的 \(\varepsilon\) 无关的常数 \(C\)，使误差绝对值不超过 \(C\varepsilon\)。它读不成``误差大约等于 \(\varepsilon\)''，也必须交代这个界对哪些 \(t\) 成立。一个近似可以在每个固定时刻都很准，却在 \(t=O(1/\varepsilon)\) 的长时间尺度上累积出 \(O(1)\) 的误差；带振荡相位时更麻烦，极小的频率偏差会慢慢把整段波形错开。

同样的双时间尺度结构在训练里也随处可见。自适应优化器的一阶、二阶矩以远快于参数本身的速度追踪梯度统计，规范化层的滑动统计量同理；强化学习里目标网络被刻意调慢，让价值估计在一个准静止的目标下收敛。这些设计能成立，靠的都是同一个前提——快变量确实来得及在慢变量动之前落到平衡。前提一旦不成立，比如两个尺度靠得太近，图上那条贴合曲线的轨迹就会变成来回震荡的螺旋，训练表现为损失曲线的持续抖动。

所以渐近分析不是删去小项。可靠的流程是：先无量纲化定出小参数，再判断令它为零是否改变方程的阶数或边界条件，然后构造与尺度相配的近似，最后交代误差随参数和观察区间如何变化。快变量迅速稳定时，慢变量模型往往够用；快层触及约束或失稳时，忽略它就漏掉了风险本身。
